\section{Isolation: what a boundary costs, and what it buys}
\label{sec:isolation}

Sandbox comparisons are usually published as one column of milliseconds, and
that column conflates two primitives that are not alternatives. We measured both
on the same host, 20 runs each.

\begin{table}[h]
\centering
\begin{tabular}{lrrl}
\toprule
\textbf{Primitive} & \textbf{median} & \textbf{p95} & \textbf{What it can run} \\
\midrule
V8 context, in-process   & \rVeight{}\,ms  & \rVeightP{}\,ms  & JavaScript. No shell, no pip, no pytest. \\
microVM, cold boot       & \rVm{}\,ms      & \rVmP{}\,ms      & Any language, any binary, a real kernel and filesystem. \\
microVM, from a checkpoint & \rVmCkpt{}\,ms & \rVmCkptP{}\,ms & As above, on the disk an earlier run left behind. \\
\bottomrule
\end{tabular}
\caption{Cold start by isolation primitive, 20 runs each. The measured ratio
between a microVM boot and an in-process V8 context is
\rVmRatio{}$\times$.}
\label{tab:sandbox}
\end{table}

\paragraph{The ratio is real and it is not a defect.}
\rVmRatio{}$\times$ is what a kernel-enforced boundary costs: a kernel to boot,
a root image to attach, a virtual network interface, and a machine to shut down. A V8 context
is a data structure in a process that is already running, so it starts in
\rVeight{}\,ms and offers a boundary that is exactly as strong as the embedding
process's correctness. These buy different things and the choice is per workload,
not per millisecond.

The workload decides it, and for agents the workload is frequently unambiguous.
An agent asked to run a project's test suite needs a package manager, a
filesystem and a shell. No JavaScript isolate provides those at any latency. The
right comparison for that workload is the second row against other people's
second rows, and \emph{we do not win it}. \rVm{}\,ms is slower than the
under-200\,ms cold start published for Firecracker-based sandboxes, and faster
than the container-based platforms whose published figures run to
seconds~\cite{e2b,modal,cfworkers}. Those figures were taken by their authors on
their own hardware under protocols we did not run; that is the reason none of
these numbers settles the question, and it is not an argument that ours is
better than it measures.

\paragraph{A checkpoint saves the disk, not the boot, and there is no pooled
tier.}
An earlier version of this section reported a pooled row at $37.5$\,ms and argued
that it, rather than the cold number, was what a fleet experiences. That row was
a command dispatched into an already-running OCI container, measured on a runtime
this system does not use. The runtime it does use offers no such mode: a VM is
booted, a command runs, the machine stops. Starting \emph{from a checkpoint} is
the nearest thing and it is not the same thing --- \rVmCkpt{}\,ms against
\rVm{}\,ms for a cold boot, which is no difference, because a checkpoint carries
the filesystem a previous run left and the kernel boots either way.

The published resume figures for other microVM platforms are memory snapshots,
which is a different mechanism~\cite{e2b}. Setting our checkpoint number beside
theirs would claim a capability by borrowing the name of one. The honest position
is that every call here pays a boot, and that the cost of \emph{not} having a
warm tier is the difference between the two rows of Table~\ref{tab:sandbox} and
the first row: roughly $3\times$ on a whole agent call once search and embedding
are counted with it.

\paragraph{Three tiers, chosen by threat model.}
Taken with \S\ref{sec:code}, the measurements describe a ladder, and the sensible
policy is to take the cheapest rung the threat model permits:

\begin{enumerate}[leftmargin=*]
\item \textbf{In-process call} --- our own code, no boundary, nanoseconds.
\item \textbf{In-process runtime} --- tenant code that is not adversarial, a
  capability boundary, \rWazero{}--\rGpython{}\,\textmu s (\S\ref{sec:code}).
\item \textbf{microVM} --- arbitrary binaries, an adversarial tenant, or a real
  filesystem: \rVm{}\,ms, on every call, because there is no warm tier.
\end{enumerate}

A platform that offers only the third rung charges every workload the price of
the most demanding one. That is the cost this paper is trying to make visible ---
not the container itself, which is worth its price when it is needed.
